% Agentic Fraud Detection System Paper
% Based on Springer LNCS template

\documentclass{svproc}

\usepackage{url}
\usepackage{helvet}
\usepackage{courier}
\usepackage{color}
\usepackage{makeidx}
\usepackage{graphicx}
\usepackage[misc]{ifsym}
\usepackage{multicol}
\usepackage[bottom]{footmisc}
\usepackage{tabularx}
\usepackage{longtable}
\usepackage{rotating}
\usepackage{amsmath}
\usepackage{subfig}
\usepackage{algorithm}
\usepackage{algorithmic}

\def\UrlFont{\rmfamily}

\begin{document}
\mainmatter

\title{Automated Fraud Detection System using MAPE-K Architecture}

\titlerunning{Automated Fraud Detection using MAPE-K}  

\author{Priti Gupta\inst{1}}

\authorrunning{P. Gupta}

\institute{Your Institution, City, Country \\
\email{youremail@institution.com}}

\maketitle

\begin{abstract}
Fraud detection systems face significant challenges including class imbalance, data drift, and adversarial evasion attacks. This paper presents an automated fraud detection system built on the MAPE-K (Monitor-Analyze-Plan-Execute using Knowledge) architecture that enables continuous autonomous adaptation after deployment. The system integrates: social media-based data collection and LLM-driven annotation, CTGAN-based data augmentation, adversarial training for robustness, and a closed-loop MAPE-K framework that operates continuously after initial deployment. Our system achieves F1 of 0.9911, ROC-AUC of 0.9966, and 98.34\% robustness. The key contribution is the autonomous capability: after deployment, the system continuously monitors for drift, automatically triggers retraining when needed, and uses a Knowledge Base to learn from past adaptation cycles. The MAPE-K loop detects drift via PSI/KS, triggers automatic retraining when needed, and includes a Supervisor with convergence monitoring. Unlike static pipelines, our system adapts and improves over time through continuous autonomous adaptation.

\keywords{Fraud Detection, MAPE-K Architecture, Autonomous Adaptation, Adversarial Training, Data Drift, CTGAN, Automated Systems}
\end{abstract}

\section{Introduction}\label{sec:intro}
Fraud is a major issue in today's digital economy. Financial losses from fraudulent transactions reach billions of dollars annually, affecting both businesses and individuals. 

When we started exploring fraud detection, we found three key challenges. First, most datasets have far more normal transactions than fraudulent ones, making it hard for models to learn what fraud looks like. Second, fraud patterns change over time as scammers develop new techniques - this is called data drift. Third, we realized that fraudsters can deliberately manipulate their behavior to bypass detection systems.

Most existing approaches use single machine learning models that cannot adapt when patterns change. Some systems use rule-based methods, but these struggle to keep up with evolving fraud. Other solutions exist for individual problems like class imbalance or adversarial attacks, but fewer systems handle all three challenges together.

This paper aims to build an automated fraud detection system that can adapt automatically. We use the MAPE-K architecture, where automated modules monitor the system, analyze conditions, plan responses, and execute actions. This allows the system to detect when fraud patterns shift and retrain itself without human help.

We gathered data from social media platforms where people share their scam experiences. Large language models help label this data automatically. To handle class imbalance, we use CTGAN to create more examples of normal transactions. We also train the model with adversarial examples so it stays robust even when attacked. The MAPE-K agents tie everything together, checking for drift and deciding when to retrain.

The main contributions of this work are a complete automated fraud detection pipeline that collects data, augments it, trains robustly, and adapts on its own, implementation of MAPE-K architecture for fraud detection that can measure how well the model learns over time, demonstration that combining CTGAN with adversarial training achieves 98.34\% robustness, and demonstrates promising experimental results F1 score of 0.9911 and ROC-AUC of 0.9966.

\section{Related Work}\label{sec:related}
Fraud detection has been widely studied in the literature. Traditional approaches use rule-based systems and basic machine learning models like logistic regression and decision trees. More recent work applies deep learning methods to detect fraudulent transactions. However, these approaches typically treat fraud detection as a static classification problem and do not address the challenges of changing fraud patterns.

Class imbalance is a well-known challenge in fraud detection. Dal Pozzolo et al. \cite{dalpozzolo2015credit} studied credit card fraud detection and showed that class imbalance significantly affects model performance. Various techniques have been proposed to handle this, including undersampling, oversampling, and synthetic data generation. CTGAN, introduced by Xu et al. \cite{xu2019ctgan}, generates synthetic tabular data that preserves the statistical properties of the original dataset. Several works have applied CTGAN to fraud detection with promising results \cite{strelcenia2023synthetic}.

Concept drift is another critical issue in fraud detection. As fraudsters change their tactics, the data distribution shifts over time. Goudar et al. \cite{goudar2026adaptive} proposed adaptive bagging strategies to address data drift in fraud detection systems. Various drift detection methods exist, including population stability index and Kolmogorov-Smirnov tests.

Adversarial robustness has gained attention in recent years. Madry et al. \cite{madry2018adversarial} showed that adversarial training can improve model robustness against evasion attacks. For tabular data, Simonetto et al. \cite{simonetto2024tabularbench} introduced TabularBench to benchmark adversarial robustness. Melo et al. \cite{melo2023adversarial} explored adversarial training specifically for tabular fraud detection data.

The MAPE-K architecture, introduced by IBM \cite{ibm2005autonomic}, provides a framework for building self-managing systems. Recent work has applied MAPE-K to various domains including microservices \cite{bucchiarone2022mape} and cloud systems \cite{oh2023mape}. However, limited work has explored MAPE-K for fraud detection specifically.

Our work differs from existing approaches by building an automated system that integrates data collection, augmentation, adversarial training, and autonomous adaptation using the MAPE-K architecture.

Recent work has explored adaptive fraud detection approaches. Chen et al. \cite{chen2024adaptive} proposed a drift-aware ensemble that updates model weights based on PSI metrics, while Rodriguez et al. \cite{rodriguez2023continual} used continual learning with Elastic Weight Consolidation. However, these approaches focus narrowly on drift adaptation without integrating the full lifecycle of data collection, augmentation, and adversarial robustness. Our work uniquely combines these elements within the comprehensive MAPE-K framework for end-to-end autonomous fraud detection.

\section{Methodology}\label{sec:methodology}
This section describes the four main components of our automated fraud detection system. Our approach was to tackle each challenge systematically: first handle the imbalanced data, then make the model robust against adversarial attacks, and finally ensure the system can adapt automatically when fraud patterns change.

\subsection{Data Collection and Annotation}
We collected fraud-related data from social media platforms where people share their scam experiences. When we looked at available datasets, we realized traditional sources like credit card data were either private or had limitations. Social media emerged as a rich source where victims share their experiences openly. The data collection process involved two main steps.

First, we scraped posts from Reddit fraud-focused communities (e.g., r/Scams, r/antiScam, r/fraud) along with their comments, using targeted keywords related to common fraud types such as investment scams, phishing, and payment fraud. The scraping targeted posts from the past several months to capture recent fraud patterns. We collected post content, comments, timestamps, and user interactions.

Second, we used the Llama3 Instruct model to automatically label the collected data. Posts and comments were transformed into numerical features using TF-IDF vectorization with n-gram range (1,2) and maximum features=5000. Comments were aggregated per post using mean pooling of TF-IDF vectors. Text was lowercased, URLs replaced with '<URL>', and user mentions replaced with '<USER>'. The LLM was prompted as an expert fraud analyst to identify whether each post described actual fraud using the post and its comments, extract the type of fraud if present, note the currency involved, extract the amount mentioned, and identify psychological tactics used (e.g., urgency, fear, authority). The LLM provided three labels: fraud (1), non-fraud (0), and uncertain (-1), along with confidence scores. Posts with confidence scores below 0.7 or labeled as uncertain (-1) were filtered out to ensure data quality. This two-stage filtering (uncertain labels and low-confidence) ensures high-quality labeled examples for training.

To validate the Llama3 annotation quality, we conducted a multi-stage validation. First, 200+ samples were verified using GPT-5 to cross-check Llama3 annotations. Second, 100+ samples were manually verified by the authors. This confirms the reliability of the Llama3-based annotation approach.

\subsection{Data Augmentation using CTGAN}
Many fraud detection datasets have class imbalance, with fraud as the minority class. However, our Reddit-sourced dataset has the opposite distribution: approximately 95% fraud and 5% non-fraud, because we specifically scraped from fraud-focused communities where users share their scam experiences. This makes non-fraud the minority class at only 5%. We use CTGAN, a generative adversarial network designed for tabular data \cite{xu2019ctgan}, to generate synthetic non-fraud examples to balance the dataset.

CTGAN consists of two neural networks that work against each other. The generator creates synthetic data that looks real, while the discriminator tries to distinguish between real and fake data. Through training, the generator learns to produce realistic samples that match the statistical properties of the original data.

We trained CTGAN on our minority class examples, which in this case are the normal transactions. The trained model generates new synthetic examples of normal transactions. These synthetic examples help balance the dataset, giving the model more examples to learn from.

To ensure the quality of synthetic data, we validated it using statistical tests. We compared the distribution of synthetic data with real data using Jensen-Shannon Divergence to ensure the generated samples maintain the same statistical properties as the original dataset.

\subsection{Adversarial Training}
Fraudsters can deliberately try to trick detection systems by modifying their transaction patterns. To make our model robust against such attacks, we employ adversarial training using the Fast Gradient Sign Method (FGSM) \cite{madry2018adversarial}. FGSM is particularly suitable for tabular fraud detection data as it generates computationally efficient adversarial perturbations that effectively simulate realistic evasion attempts while requiring minimal domain-specific knowledge about feature interactions.

Given an input feature vector $\mathbf{x}$ with label $y$, the adversarial example is generated as:

\begin{equation}
\mathbf{x}_{adv} = \mathbf{x} + \epsilon \cdot \text{sign}(\nabla_{\mathbf{x}} \mathcal{L}(\theta, \mathbf{x}, y))
\end{equation}

where $\epsilon = 0.05$ is the perturbation magnitude (5\% of feature range), $\nabla_{\mathbf{x}} \mathcal{L}$ is the gradient of the loss function with respect to the input features, and $\theta$ represents the model parameters. The sign operation ensures perturbations are in the direction that maximizes the loss.

During training, the model is exposed to both clean examples $(\mathbf{x}, y)$ and adversarial examples $(\mathbf{x}_{adv}, y)$. This enables the model to learn robust decision boundaries that remain effective even when input features are maliciously perturbed. The model is trained to correctly classify both clean and adversarial examples, improving robustness against evasion attacks without significantly degrading performance on clean data \cite{goodfellow2014gan}.

We evaluate robustness by measuring performance degradation under different perturbation strengths $\epsilon \in \{0.01, 0.05, 0.10, 0.20, 0.30\}$. The robustness score is calculated as the ratio of F1 score under attack to clean F1 score.

\subsection{Autonomous Adaptation Layer: MAPE-K Architecture}

The MAPE-K architecture consists of five specialized agents that work together to enable autonomous adaptation. The Ingestion Agent collects new data from social media sources and preprocesses it for analysis. The Drift Agent continuously monitors incoming data using Population Stability Index (PSI) and Kolmogorov-Smirnov (KS) tests to detect data drift by comparing current distributions against historical baselines stored in the Knowledge Base. The Evaluation Agent assesses model performance using F1-score, ROC-AUC, and robustness metrics, comparing against baseline performance to determine if adaptation is needed. The Planning phase (led by the Training Agent) defines the adaptation strategy including retraining parameters, and includes the Decision Engine which determines whether to trigger retraining based on Analysis outputs. The Execution phase (led by the Balance Agent) handles data rebalancing using CTGAN and manages model deployment.

Specifically, the Decision Engine (part of Planning) determines the operational path based on outputs from the Drift and Evaluation Agents. If performance meets thresholds and no drift is detected, the current model continues in production. Otherwise, corrective adaptation is triggered.

The Knowledge Base serves as a centralized historical repository that all agents access and update. It stores drift statistics, performance metrics, training configurations, and deployment history, enabling cumulative learning over time.

The critical distinction is that the MAPE-K loop operates continuously after deployment, not during initial training. Each adaptation cycle enables the system to learn from previous experiences, making it a genuine autonomous system.

The system thus achieves autonomous capability by detecting degradation triggers (drift or poor performance), planning appropriate responses based on historical knowledge (data or model-level corrections), executing changes, and verifying outcomes---all while leveraging the Knowledge Base to improve future decisions through accumulated experience.

\section{Experimental Setup}\label{sec:experimental}
This section describes our experimental setup including the dataset, model configuration, and evaluation metrics.

\subsection{Dataset}
We used a dataset collected from social media platforms containing labeled fraud examples. The dataset consists of posts labeled as either fraudulent or normal (non-fraud). The data was collected over a period of several months from Reddit fraud-focused communities.

The raw dataset contained 9,000 posts scraped from Reddit fraud forums. The LLM annotation provided three labels: fraud (1), non-fraud (0), and uncertain (-1). After removing approximately 4,000 posts labeled as uncertain, we retained 5,000 labeled examples. For experimentation, we used a subset of 3,000 examples.

Unlike typical fraud datasets where fraud is the minority class (e.g., credit card fraud), our dataset contains more fraud examples than non-fraud examples because we specifically scraped from fraud-focused communities. Our data distribution is approximately 95% fraud and 5% non-fraud. Therefore, CTGAN generates non-fraud (the actual minority class) to balance the dataset, which is the correct approach for our data distribution.

For data drift experiments, we simulated drift by splitting our dataset chronologically, using data from earlier months for training and later months for testing, thereby introducing natural temporal variations in fraud tactics. Additionally, for controlled experiments, we injected drift by modifying 20% of non-fraud samples to mimic emerging fraud patterns (e.g., changing currency references, urgency tactics).

\subsection{Model Configuration}
We used XGBoost as our base classifier. XGBoost is a gradient boosting algorithm that works well with tabular data and has been successfully applied to fraud detection tasks. The model was trained with the following key parameters.

The learning rate was set to 0.1, which controls how much the model updates its weights with each iteration. The maximum depth was limited to 6 to prevent overfitting. The number of estimators was set to 100. We used binary cross-entropy as the objective function.

For adversarial training \cite{madry2018adversarial}, we employed FGSM with perturbation magnitude $\epsilon = 0.05$ (5 percent of feature range). The model was trained on both clean and adversarial examples to learn robust classification boundaries.

The CTGAN \cite{xu2019ctgan} was trained with 100 epochs and a batch size of 16. The generator and discriminator dimensions were set to 64. These values were chosen based on preliminary experiments to balance training time and quality.

To evaluate robustness, we used the TabularBench framework \cite{simonetto2024tabularbench} which provides standard benchmarks for adversarial robustness on tabular data.

\subsection{Evaluation Metrics}
We evaluated our system using several standard metrics.

For classification performance, we used accuracy, precision, recall, and F1 score. Accuracy measures the overall proportion of correct predictions. Precision measures how many predicted fraud cases are actually fraud. Recall measures how many actual fraud cases are detected. F1 score is the harmonic mean of precision and recall, providing a balanced measure.

We also used ROC-AUC, which measures the area under the receiver operating characteristic curve. This metric evaluates the model's ability to distinguish between fraud and normal cases across different decision thresholds.

For drift detection \cite{goudar2026adaptive}, we used Population Stability Index and Kolmogorov-Smirnov test. PSI values above 0.20 indicate significant drift requiring action. KS values above 0.05 indicate significant distribution changes.

For robustness evaluation, we measured the model's performance under adversarial perturbations at different epsilon values. The robustness score is calculated as the ratio of F1 under attack to clean F1. We target a robustness score above 0.60.

For data quality, we used Jensen-Shannon Divergence to compare synthetic and real data distributions. JSD values below 0.20 indicate high-quality synthetic data that maintains statistical properties of the original dataset.

\subsection{Experimental Environment}
All experiments were conducted on a standard workstation with Python implementation. Key libraries used include scikit-learn 1.3.0 for model training and evaluation, XGBoost 1.7.0 for gradient boosting, and SDV 0.14.0 for CTGAN. To ensure reproducibility, all experiments were conducted with random seeds fixed to 42. The system was run multiple times to ensure consistency of results.

\section{Results}\label{sec:results}
This section presents the experimental results for our automated fraud detection system. We evaluated performance across multiple dimensions including classification accuracy, robustness, drift detection, and data quality.

\subsection{Classification Performance}
Our system achieves promising results on the fraud detection task. The model reached an F1 score of 0.9911, indicating reasonable precision and recall balance. The ROC-AUC score of 0.9966 shows reasonable ability to distinguish between fraud and normal transactions across different thresholds.

The precision and recall both reached 0.9911, meaning the model correctly identifies fraudulent transactions while keeping false alarms low. The false positive rate was 7.79 percent, which is acceptable for practical deployment where catching fraud is more important than slightly annoying users with occasional false alarms.

These results demonstrate that combining CTGAN data augmentation with adversarial training produces an effective fraud detector that performs well on real-world data.

\subsection{Robustness under Adversarial Attacks}
We tested how well the model handles adversarial perturbations. We evaluated performance at different attack strengths, ranging from zero to 0.30 perturbation levels.

The model maintained strong performance even under attack. At the highest perturbation level of 0.30, the F1 score remained at 0.9803, only dropping slightly from the clean F1 of 0.9911. This shows the model is robust against attempts to trick it by modifying transaction features.

The overall robustness score was 98.34 percent, calculated as the ratio of average F1 under attack to clean F1. This far exceeds our target threshold of 60 percent. The attack success rate, measuring how much performance degrades under attack, was only 1.09 percent.

These results demonstrate that adversarial training effectively prepares the model to handle evasion attempts. Even when fraudsters try to manipulate their transactions to bypass detection, the model maintains high accuracy.

\subsection{Drift Detection and Adaptation}
We tested the system's ability to detect and adapt to data drift. During one of our experimental runs, significant drift was detected with a PSI value of 0.22, which exceeded our threshold of 0.20. The KS statistic also showed drift at 0.09.

When drift was detected, the system automatically triggered the correction cycle through the MAPE-K loop: Monitor detected drift, Analyze assessed metrics, Decision routed to Plan (Training Agent), then Execute (Balance Agent + Deploy).

After retraining, the system successfully recovered. The new model achieved F1 score of 0.92, showing the system can adapt to changing fraud patterns. The drift metrics returned to normal levels, confirming the adaptation was successful.

This demonstrates the practical value of the MAPE-K architecture. The system can detect problems on its own and take corrective action without human intervention.

\subsection{Data Quality Validation}
We validated the quality of synthetic data generated by CTGAN using Jensen-Shannon Divergence. The JSD score was 0.1291, which is below our threshold of 0.20. This confirms the synthetic data maintains the statistical properties of the original dataset.

We also performed chi-squared tests on categorical features and Kolmogorov-Smirnov tests on numerical features. The categorical test passed for 100 percent of features. The numerical test passed for 97.1 percent of features. These results confirm the synthetic data is realistic and suitable for training.

To find the optimal amount of synthetic data, we experimented with different quantities. The best results were achieved with around 700 synthetic samples. Adding more than this did not improve performance and in some cases hurt it due to noise introduction.

\subsection{Comparison with Baseline and Ablation Study}
We conducted both baseline comparison and ablation studies to understand the contribution of each component.

\begin{table}[h]
\centering
\caption{Ablation Study: Component-wise Performance on Clean Data and Under Drift Conditions}
\label{tab:ablation}
\begin{tabular}{|l|c|c|c|c|}
\hline
\textbf{Configuration} & \textbf{F1 (clean)} & \textbf{F1 (drift)} & \textbf{ROC-AUC} & \textbf{Robustness (\%)} \\
\hline
Base (XGBoost only) & 0.74 & 0.74 & 0.78 & 45.2 \\
+ CTGAN & 0.85 & 0.87 & 0.91 & 48.7 \\
+ Adversarial Training & 0.91 & 0.92 & 0.95 & 89.3 \\
+ Full MAPE-K Loop & 0.99 & 0.92 & 1.00 & 98.34 \\
\hline
\end{tabular}
\end{table}

The ablation study shows component contribution under both clean and drift conditions. Under clean data, the base model achieves F1 of 0.74, CTGAN improves it to 0.85, and the full MAPE-K loop achieves 0.99. Under drift, the full MAPE-K loop recovers to 0.92, demonstrating autonomous adaptation. Robustness improves from 45.2% to 98.34%.

We also compared against a baseline without MAPE-K under drift conditions. A baseline model without CTGAN or adversarial training achieved F1 of 0.74 under drift conditions. Adding CTGAN improved this to 0.87. Adding adversarial training on top of CTGAN achieved F1 of 0.92 even under drift. The MAPE-K architecture provides additional value by enabling automatic adaptation - when drift occurs, the baseline without MAPE-K would require manual intervention to retrain. Our system handles this automatically.

The combination of all components produces good results and enables an autonomous fraud detection system.

\section{Discussion}\label{sec:discussion}
This section discusses what our results mean and how they compare to existing work.

Our results show that combining CTGAN with adversarial training produces strong fraud detection performance. The CTGAN helps the model learn better by providing more balanced training data. Adversarial training then makes the model robust against attacks. Together, they achieve 98.34 percent robustness, meaning the model keeps working well even when someone tries to trick it.

It is important to note that these results were obtained on our specific Reddit-sourced dataset with a 95% fraud distribution. Typical fraud detection datasets (e.g., credit card fraud) have the opposite distribution with fraud as the minority class. Further validation on external datasets such as IEEE-CIS or ULB credit card fraud datasets would strengthen these findings. Results may vary on datasets with different class distributions or from different sources.

The MAPE-K architecture adds real value by enabling automatic adaptation. Traditional fraud detection systems need humans to notice when something is wrong and manually retrain the model. Our system handles this automatically. When drift is detected, the system retrains itself and checks if performance improves. This reduces the need for human oversight and ensures the system stays effective as fraud patterns change.

One interesting finding is that the optimal amount of synthetic data is around 700 samples. Adding more did not help and sometimes hurt performance. This suggests that while synthetic data is useful, too much can introduce noise. Finding the right balance is important.

Our system performed well even under data drift. When fraud patterns changed, the model could still maintain F1 above 0.90 after automatic retraining. This is better than many static systems that would fail completely under the same conditions.

There are some limitations to consider. First, we collected data from social media which may not represent all types of fraud. Some fraud schemes may not be discussed publicly online. Second, we used XGBoost which is a traditional machine learning model. More complex deep learning approaches might achieve even better results. Third, while we use FGSM-based adversarial training, more sophisticated attacks like PGD could be explored in future work to further test robustness.

The system demonstrates strong performance metrics that suggest potential for real-world application. The F1 score of 0.9911 and robustness of 98.34 percent are encouraging, though further validation would be needed for production deployment. The automatic adaptation means the system can handle changing fraud patterns without constant human attention.

Compared to related work, our system integrates multiple components that others typically handle separately. Most papers address class imbalance or adversarial training in isolation. Our complete pipeline from data collection to autonomous adaptation integrates these components in a unified framework.

Future work could explore different model architectures, test on more diverse datasets, and implement more advanced adversarial training methods. The MAPE-K framework could also be extended with more sophisticated planning algorithms.

Overall, our automated fraud detection system demonstrates how combining established techniques like CTGAN and adversarial training with the MAPE-K architecture creates a robust and adaptive solution for real-world fraud detection challenges.

\section{Conclusion}\label{sec:conclusion}
This paper presented an automated fraud detection system built on the MAPE-K architecture. The system integrates five specialized agents (Ingestion, Drift, Evaluation, Training, Balance) following the M $\rightarrow$ A $\rightarrow$ P $\rightarrow$ E order, where the Decision Engine is part of the Planning phase, with feedback loops and Supervisor oversight.

Our experiments showed promising results. The system achieves F1 score of 0.9911 and ROC-AUC of 0.9966 on real fraud data. The robustness against adversarial attacks is 98.34 percent, demonstrating that the model can handle attempts to evade detection. The MAPE-K architecture successfully detects data drift and triggers automatic retraining, achieving F1 of 0.92 even under changing fraud patterns.

The key contributions of this work are a complete automated fraud detection pipeline that handles data collection, augmentation, adversarial training, and autonomous adaptation, a practical implementation of MAPE-K for fraud detection that demonstrates measurable learning adaptation, evidence that combining CTGAN with adversarial training produces robust models, and strong experimental results validated on real-world data.

For practitioners, this system offers a deployable solution that reduces manual oversight through automatic adaptation. For researchers, the integration of multiple techniques into a coherent framework provides a foundation for further work.

Future directions include exploring more advanced model architectures, testing on additional datasets from different sources, and implementing more sophisticated adversarial training methods. The MAPE-K framework could also be extended with better planning algorithms to optimize retraining decisions.

In conclusion, our automated fraud detection system demonstrates how automated modules can work together to handle the complex challenges of modern fraud detection. By combining established techniques with autonomous adaptation, we create a robust and practical solution for real-world fraud detection needs.

\bibliographystyle{splncs04}
\bibliography{references}

\end{document}